% PAGES: 291-292
% Chunk c3 (pages 291-302, folios 275-286), file 1 of 5.
% ch09/ was completely empty when this chunk started (re-listed immediately
% before beginning this file); chunks c1 (267-278) and c2 (279-290) had not
% yet written anything at that point. Both had finished by the time this
% chunk was ready to submit (re-listed again per the standing brief), and
% c2's last file, sec09_c2_12.tex, explicitly confirms a clean close: "The
% file ends cleanly at a paragraph boundary (...is $1-C$.) with no open
% \begin{...}, no open \[, and no open non-environment example block". That
% matches what page 291 (folio 275) itself shows: a fresh paragraph ("A
% frequent assumption is that the return of a portfolio is normally
% distributed.") that is not mid-sentence and does not continue an open
% display or environment. Nothing needs to be closed here. This file
% continues inside section 9.6 "Value at Risk (VaR). Portfolio VaR." (opened
% by c2), referencing c2's formulas (9.84) and (9.85) by number.

A frequent assumption is that the return of a portfolio is normally
distributed. Note that this would indeed be the case if the returns of the
portfolio components have a joint multivariate normal distribution. Then, any
linear combination of their returns, including the weighted average of the
portfolio returns that gives the portfolio return, see (9.2), would be
normally distributed; cf. Theorem 7.7.

If a portfolio is assumed to have normally distributed returns, then the
formulas (9.84) and (9.85) for the $N$ day $C\%$ VaR of the portfolio hold. If
$\mu_{R_V}$ and $\sigma_{R_V}$ are the annualized expected return and the
standard deviation of the return of the portfolio over a small time period
$t$, then
\begin{equation}
\frac{V(t)-V(0)}{V(0)} \;=\; \mu_{R_V} t \;+\; \sigma_{R_V}\sqrt{t}\,Z.
\end{equation}

For an $N$ day horizon, let $t = \frac{N}{252}$ in (9.86) and obtain that
\begin{equation}
\frac{V(N)-V(0)}{V(0)} \;=\; \mu_{R_V}\frac{N}{252} \;+\; \sigma_{R_V}\sqrt{\frac{N}{252}}\,Z,
\end{equation}
where $V(N)$ is a simplified notation for $V\left(\frac{N}{252}\right)$, the
value of the portfolio in $N$ days.

From (9.85) and (9.87), it follows that\footnote{The assumption that the
portfolio has positive initial value, i.e., $V(0) > 0$, was made and used
implicitly in (9.88).}
\begin{align}
&P\bigl(V(N) - V(0) \;\le\; -\VaR(N,C)\bigr) \;=\; 1 - C \notag\\
\Longleftrightarrow\quad &P\left(\frac{V(N)-V(0)}{V(0)} \;\le\; -\frac{\VaR(N,C)}{V(0)}\right) \;=\; 1-C \\
\Longleftrightarrow\quad &P\left(\mu_{R_V}\frac{N}{252} + \sigma_{R_V}\sqrt{\frac{N}{252}}\,Z \;\le\; -\frac{\VaR(N,C)}{V(0)}\right) \;=\; 1-C \notag\\
\Longleftrightarrow\quad &P\left(\sigma_{R_V}\sqrt{\frac{N}{252}}\,Z \;\le\; -\frac{\VaR(N,C)}{V(0)} - \mu_{R_V}\frac{N}{252}\right) \;=\; 1-C \notag\\
\Longleftrightarrow\quad &P\left(Z \;\le\; -\sqrt{\frac{252}{N}}\,\frac{\VaR(N,C)}{\sigma_{R_V}V(0)} - \sqrt{\frac{N}{252}}\,\frac{\mu_{R_V}}{\sigma_{R_V}}\right) \;=\; 1-C \notag\\
\Longleftrightarrow\quad &P\left(Z \;\le\; -\left(\sqrt{\frac{252}{N}}\,\frac{\VaR(N,C)}{\sigma_{R_V}V(0)} + \sqrt{\frac{N}{252}}\,\frac{\mu_{R_V}}{\sigma_{R_V}}\right)\right) \;=\; 1-C.
\end{align}

Let $z_C$ be the $z$--score of the standard normal distribution corresponding
to $C$, i.e., $P(Z \le z_C) = C$. Note that
\[
\begin{aligned}
z_{99} &= 2.3263, & \text{since}\quad & P(Z \le 2.3263) = 0.99 = 99\%\\
z_{98} &= 2.0537, & \text{since}\quad & P(Z \le 2.0537) = 0.98 = 98\%;\\
z_{95} &= 1.6449, & \text{since}\quad & P(Z \le 1.6449) = 0.95 = 95\%;\\
z_{90} &= 1.2816, & \text{since}\quad & P(Z \le 1.2816) = 0.90 = 90\%.
\end{aligned}
\]

Recall that $P(Z \le -a) = 1 - P(Z\le a)$ for any $a\in\R$; see Lemma 10.19 for
a proof of this result. Then,
\begin{equation}
P(Z \le -z_C) \;=\; 1-C.
\end{equation}

From (9.89) and (9.90), we obtain that
\begin{equation}
z_C \;=\; \sqrt{\frac{252}{N}}\,\frac{\VaR(N,C)}{\sigma_{R_V}V(0)} \;+\; \sqrt{\frac{N}{252}}\,\frac{\mu_{R_V}}{\sigma_{R_V}}.
\end{equation}

By solving (9.91) for $\VaR(N,C)$, we find that
\begin{equation}
\VaR(N,C) \;=\; \sqrt{\frac{N}{252}}\,\sigma_{R_V} z_C V(0) \;-\; \frac{N}{252}\mu_{R_V}V(0).
\end{equation}

For small values of $N$, e.g., for $N \le 10$ corresponding to time horizons
of less than two weeks, $\frac{N}{252}$ is much smaller\footnote{For example,
for $N=10$, $\sqrt{\frac{N}{252}} = \sqrt{\frac{10}{252}} \approx 0.20$ and
$\frac{N}{252} = \frac{10}{252} \approx 0.04$.} than $\sqrt{\frac{N}{252}}$.
Thus, from (9.92), we obtain the following approximation:
\begin{equation}
\VaR(N,C) \;\approx\; \sqrt{\frac{N}{252}}\,\sigma_{R_V} z_C V(0).
\end{equation}

From (9.93), it can be inferred that the Value at Risk of a portfolio with
normally distributed return is proportional to the square root of the time
horizon, which provides the following simple way to scale the Value at Risk
of a portfolio to different time horizons and confidence levels:
\[
\frac{\VaR(N_2,C_2)}{\VaR(N_1,C_1)} \;\approx\; \frac{z_{C_2}\sqrt{N_2}}{z_{C_1}\sqrt{N_1}},
\]
which can also be written as
\begin{equation}
\VaR(N_2,C_2) \;\approx\; \frac{z_{C_2}\sqrt{N_2}}{z_{C_1}\sqrt{N_1}}\,\VaR(N_1,C_1).
\end{equation}

\par\medskip\noindent\textit{Example:}\ What is the ten--day 99\% VaR of a
portfolio with a five--day 98\% VaR of \$10 million?

\begin{answerx}
Assuming that the returns of the portfolio are normally distributed, we
obtain from (9.94) that
\[
\begin{aligned}
\VaR(10\text{ days},99\%) &\approx \frac{z_{99}\sqrt{10}}{z_{98}\sqrt{5}}\,\VaR(5\text{ days},98\%) \;=\; \frac{2.3263\sqrt{10}}{2.0537\sqrt{5}}\cdot\$10{,}000{,}000\\
&= \$16{,}019{,}307.
\end{aligned}
\]
\hfill$\square$
\end{answerx}

\par\medskip\noindent\textit{Example:}\ Consider two assets with multivariate
normal distribution of their returns. Their returns have standard deviations
of 25\% and 35\%, and a correlation of -25\%.

(i) Use formula (9.93) to find approximate values for the 10--day 98\% VaR of
a \$100 million portfolio if the portfolio is invested in the first asset,
invested in the second asset, or invested in the minimum variance portfolio
fully invested in the two assets.

(ii) Assume that the returns of the two assets have expected values of 3\%
and 6\%, respectively. Use formula (9.92) to find approximate values for the
10--day 98\% VaR of a \$100 million portfolio if the portfolio is invested in
the first asset, invested in the
% --- END OF FILE: mid-sentence, no environment open ---
% Page 292 (folio 276) ends here mid-sentence, in the middle of part (ii) of
% this Example's plain-text (not \begin{examplex}) question. No LaTeX
% environment is open at this point (the manual "Example:" run-in above is
% plain prose, not an environment, per the established book convention --
% see e.g. ch04/sec04_c2_1.tex -- since the source's "Example:" here has no
% closing square of its own; only the "Answer:" that follows on page 293
% closes with one, via \answerx). sec09_c3_2.tex continues this exact
% sentence verbatim ("second asset, or invested in the minimum variance
% portfolio...") and finishes part (ii) and (iii) as plain text, followed by
% \begin{answerx}...\end{answerx}.
